% ============================================================================
% Fichier  : partie2/P02_C05_mission_principes_cycle.tex
% Titre    : La Mission d'Audit : Principes et Cycle
% Auteur   : MINKA MI NGUIDJOÏ Thierry Emmanuel
% Version  : 1.0 -- octobre 2026
% ============================================================================
% Sources : ISO 19011:2018 (sept principes ; l'édition 2011 du cours Audit_SI
% en comptait six) ; Normes mondiales d'audit interne de l'IIA (2024),
% domaine III (charte, rattachement au conseil) ; boîte à outils O0 à O9 du
% cours Audit_SI de l'auteur, mise en cohérence avec les articles A0 et A1
% (gabarit canonique, OUI/NON/NEK, µ, U, H).
% ============================================================================

\chapter{La Mission d'Audit : Principes et Cycle}
\label{chap:mission}

\epigraph{%
    « Une mission d'audit se gagne avant le terrain\\
    et se perd après le rapport. »%
}{Adage des chefs de mission}

\niveauChapitre{Fondamental}{%
    Maîtrise du Chapitre~\ref{chap:audit-auditologie} (audit, assurance,
    axiomes de l'auditologie)%
}

\begin{boxobjectifs}
\begin{itemize}
    \item Énoncer les sept principes de l'audit selon l'ISO~19011:2018.
    \item Expliquer le rôle de la charte d'audit et de la lettre de mission.
    \item Décrire les cinq phases de la mission, leurs entrées et leurs
          livrables.
    \item Situer les dix outils de la méthode (O0 à O9) dans le cycle.
    \item Organiser le contradictoire et le relier à l'opposabilité des
          conclusions (axiome~6).
    \item Montrer comment le contradictoire réduit l'incertitude de la
          mission.
\end{itemize}
\end{boxobjectifs}

\bigskip

% ============================================================================
\section*{Question directrice : comment ne pas refaire la première mission~?}
% ============================================================================

La première mission de LiZbiZ (Chapitre~\ref{chap:prologue}) a échoué
faute de mandat, de critères et de preuves. Pour la seconde, le cabinet
AUDIX CONSEIL n'improvise rien~: M.~BEYALA ouvre la mission par une réunion
avec le directeur général, remet une lettre de mission, annonce un
calendrier en cinq phases et prévient que chaque constat sera discuté avec
les audités avant d'être écrit.

Ce chapitre présente cette organisation. Il pose les principes qui guident
l'auditeur (section~\ref{sec:c05-principes}), les documents qui fondent sa
mission (section~\ref{sec:c05-charte}), le cycle de la mission
(section~\ref{sec:c05-cycle}), les outils qui l'outillent
(section~\ref{sec:c05-outils}) et le contradictoire qui la rend opposable
(section~\ref{sec:c05-contradictoire}). Les Chapitres~\ref{chap:planifier}
à~\ref{chap:conclure} détaillent ensuite chaque phase.

\separateur

% ============================================================================
\section{Les Principes de l'Audit}
\label{sec:c05-principes}
% ============================================================================

L'ISO~19011 donne des lignes directrices pour l'audit des systèmes de
management. Ses principes valent pour tout audit, et en particulier pour
l'audit des SI~\cite{iso19011_2018}.

\begin{tableaucours}{L{3.4cm}L{4.4cm}L{4.4cm}}{%
    \tableauHeader{Principe} &
    \tableauHeader{Sens (reformulé)} &
    \tableauHeader{Chez LiZbiZ}
}
Intégrité & Honnêteté, diligence, responsabilité & Ne pas taire un constat
qui gêne la direction \\
Présentation impartiale & Rendre compte avec exactitude et fidélité &
Ne pas grossir un risque mineur, ne pas minimiser un risque majeur \\
Conscience professionnelle & Diligence et jugement proportionnés à
l'enjeu & Ne pas lancer un test de charge à l'heure de pointe \\
Confidentialité & Protéger l'information obtenue & Ne pas divulguer les
failles constatées \\
Indépendance & Impartialité, absence de conflit d'intérêts & Ne pas auditer
l'application qu'on a conçue \\
Approche fondée sur la preuve & Des conclusions vérifiables, appuyées sur
un échantillon adapté & Une capture, une extraction, pas un ouï-dire \\
Approche par les risques & Concentrer l'effort sur ce qui compte pour les
objectifs & Commencer par les risques intolérables du registre \\
\end{tableaucours}
\vspace{-0.8em}
\begin{center}{\small\itshape Tableau~5.1. Les sept principes de l'audit
selon l'ISO~19011:2018 (notre reformulation)}\end{center}

\begin{boxinfo}
\textbf{Six ou sept principes~?} L'édition de 2011 de l'ISO~19011 en
comptait six. L'édition de 2018, en vigueur, ajoute l'\emph{approche par
les risques}. Un support qui annonce « six principes » se réfère à
l'ancienne édition.
\end{boxinfo}

Deux de ces principes trouvent un fondement dans l'auditologie. L'approche
fondée sur la preuve découle de l'axiome d'opposabilité~: une conclusion
doit pouvoir être défendue devant un tiers. L'indépendance, on l'a vu au
Chapitre~\ref{chap:audit-auditologie}, se déduit de l'interaction
stratégique et de l'opposabilité~\cite{minka2026auditologie}.

% ============================================================================
\section{Charte d'Audit et Lettre de Mission}
\label{sec:c05-charte}
% ============================================================================

\subsection{La charte d'audit}

\begin{boxdef}
\termecle{Charte d'audit}\index{charte d'audit}~: document approuvé par
l'organe de gouvernance, qui fixe la mission de l'audit, ses pouvoirs (accès
aux personnes, aux documents et aux systèmes), son rattachement, son
indépendance et ses responsabilités.
\end{boxdef}

Les Normes mondiales d'audit interne de l'IIA placent sous la
responsabilité du conseil la supervision de la fonction d'audit, son mandat
et les moyens de son indépendance~\cite{iia2024gias}. La charte en est
l'instrument~: sans elle, l'auditeur dépend du bon vouloir de ceux qu'il
audite.

\subsection{La lettre de mission}

Chaque mission est ensuite fondée sur une \termecle{lettre de
mission}\index{lettre de mission}, ou des termes de référence, qui précise
ce qu'elle doit démontrer. Le Chapitre~\ref{chap:audit-auditologie} en a
donné un exemple pour LiZbiZ~; le Chapitre~\ref{chap:planifier} en détaille
le contenu.

\begin{boxalert}
\textbf{Violation typique.} Le plan d'audit interne d'une entreprise est
approuvé par le directeur général, dont relèvent les directions auditées.

\smallskip
\textbf{Règle~:} le plan d'audit et les moyens de l'audit sont approuvés par
l'organe de gouvernance (conseil, comité d'audit), et non par la direction
auditée.
\end{boxalert}

% ============================================================================
\section{Le Cycle de la Mission}
\label{sec:c05-cycle}
% ============================================================================

\subsection{Cinq phases}

\noindent\begin{minipage}{\linewidth}
\begin{center}
\begin{tikzpicture}[
    phase/.style = {
        draw=CouleurPrimaire, fill=CouleurPrimaire, rounded corners=4pt,
        text width=2.1cm, minimum height=1.1cm, align=center,
        font=\scriptsize\bfseries\color{white}
    },
    livr/.style = {
        draw=CouleurBordDef, fill=CouleurDef, rounded corners=3pt,
        text width=2.0cm, align=center, font=\scriptsize\color{CouleurPrimaire}
    },
    fleche/.style = {-{Stealth[length=5pt]}, thick, color=CouleurSecondaire}
]
\foreach \n/\t/\l [count=\i] in {
    0/Préparation/Termes de référence,
    1/Planification/Plan d'audit,
    2/Exécution/Papiers de travail,
    3/Rapport/Rapport contradictoire,
    4/Suivi/Tableau de suivi} {
    \node[phase] (p\i) at (2.55*\i, 0) {Phase~\n\\ \t};
    \node[livr, below=0.35cm of p\i] {\l};
}
\foreach \i [evaluate=\i as \j using int(\i+1)] in {1,...,4}
    \draw[fleche] (p\i) -- (p\j);
\draw[fleche] (p5.north) to[out=150, in=30] node[above, font=\scriptsize\itshape]
    {mission suivante} (p1.north);
\end{tikzpicture}
\end{center}
\vspace{-0.4em}
\begin{center}{\small\itshape Figure~5.1. Les cinq phases de la mission et
leurs livrables}\end{center}
\end{minipage}

\medskip

\begin{tableaucours}{L{2.6cm}L{4.8cm}L{4.8cm}}{%
    \tableauHeader{Phase} &
    \tableauHeader{Activités principales} &
    \tableauHeader{Chapitre}
}
0. Préparation & Mandat, termes de référence, documents sollicités,
équipe & Chapitre~\ref{chap:planifier} \\
1. Planification & Prise de connaissance, cartographie préliminaire des
risques, décomposition de l'espace d'audit, plan d'audit &
Chapitres~\ref{chap:planifier} et~\ref{chap:decomposition} \\
2. Exécution & Entretiens, tests, cotation \OUI{}~/ \NON{}~/ \NEK{}, papiers
de travail & Chapitre~\ref{chap:preuve} \\
3. Rapport & Calcul de $\mu$, $U$, $H$, observations, contradictoire,
rapport & Chapitre~\ref{chap:conclure} \\
4. Suivi & Suivi des recommandations, préparation du cycle suivant &
Chapitre~\ref{chap:conclure} \\
\end{tableaucours}
\vspace{-0.8em}
\begin{center}{\small\itshape Tableau~5.2. Les phases de la mission}\end{center}

Le cycle ne se referme pas sur lui-même~: le suivi d'une mission prépare la
suivante. C'est la dimension temporelle de l'audit, que l'auditologie pose
en axiome~: le système audité évolue, et son état à la mission suivante
dépend des interventions déclenchées par la précédente (axiome~5,
\emph{individuation temporelle}~\cite{minka2026auditologie}).

\subsection{La mission d'AUDIX CONSEIL chez LiZbiZ}

\begin{boxscenario}{Opération LIZBIZ. Le calendrier de la seconde mission}
\begin{tableaucours}{L{3.0cm}L{2.6cm}L{6.2cm}}{%
    \tableauHeader{Phase} &
    \tableauHeader{Dates (2026)} &
    \tableauHeader{Faits marquants}
}
0. Préparation & 4 -- 8 mai & Lettre de mission signée~; 21 documents
sollicités \\
1. Planification & 11 -- 22 mai & Prise de connaissance des quatre centres
de données~; décomposition du domaine D~; plan d'audit validé par le
directeur général \\
2. Exécution & 25 mai -- 26 juin & Entretiens, extractions, tests sur
quatre agences \\
3. Rapport & 29 juin -- 17 juillet & Réunion de restitution, réponses
écrites des audités, rapport final \\
4. Suivi & À partir de janvier~2027 & Revue de la mise en œuvre des
recommandations \\
\end{tableaucours}

\medskip
\textbf{Analyse} (confiance élevée). Le calendrier réserve trois semaines
au rapport et au contradictoire~: c'est le temps nécessaire pour que les
conclusions soient discutées, et donc opposables.
\end{boxscenario}

% ============================================================================
\section{La Boîte à Outils de la Méthode}
\label{sec:c05-outils}
% ============================================================================

La méthode du cours est outillée~: à chaque phase correspondent un ou
plusieurs instruments, numérotés de O0 à O9. Ils seront construits au fil
des chapitres et rassemblés en annexe.

\begin{tableaucours}{cL{3.6cm}L{2.4cm}L{5.4cm}}{%
    \tableauHeader{Outil} &
    \tableauHeader{Nom} &
    \tableauHeader{Phase} &
    \tableauHeader{Usage}
}
O0 & Aide-mémoire du cycle & Toutes & Entrées, activités et livrables de
chaque phase \\
O1 & Fiche technique de la mission & 0 & Présentation de la mission~:
objectifs, acteurs, référentiels \\
O2 & Termes de référence & 0 & Objectifs, périmètre, équipe, livrables,
calendrier \\
O3 & Liste des documents sollicités & 0 & Documents attendus, priorité,
interlocuteur~; un document essentiel absent est coté \NEK{} \\
O4 & Classeur d'aide à l'audit & 1 à 3 & Décomposition selon le gabarit
canonique, fiches d'assertion, cotation \OUI{}~/ \NON{}~/ \NEK{}, calcul de
$\mu$, $U$ et $H$ par domaine \\
O5 & Guide d'entretien & 2 & Questions issues des fiches d'assertion,
classées par interlocuteur \\
O6 & Plan d'audit & 1 & Contexte, périmètre, risques, programme de
travail, calendrier \\
O7 & Plan de tests & 2 & Tests codifiés, preuves attendues, références des
papiers de travail \\
O8 & Projets d'observation & 3 & Observations en sept éléments, soumises
au contradictoire \\
O9 & Suivi des documents & 2 & Mise à jour, sur le terrain, de la liste O3 \\
\end{tableaucours}
\vspace{-0.8em}
\begin{center}{\small\itshape Tableau~5.3. Les outils O0 à O9}\end{center}

L'outil central est le classeur~O4~: il porte la décomposition de l'espace
d'audit (Chapitre~\ref{chap:decomposition}), la cotation des réponses
(Chapitre~\ref{chap:preuve}) et les trois mesures de chaque domaine
(Chapitre~\ref{chap:conclure}).

\begin{boxmetier}{Auditeur SI}
\textbf{Ouvrir une mission}
\begin{enumerate}
    \item Vérifier que la charte ou la lettre de mission donne l'accès aux
          personnes, aux documents et aux systèmes nécessaires.
    \item Tenir une réunion d'ouverture avec le commanditaire et les
          principaux audités~: objectifs, périmètre, calendrier, règles du
          contradictoire.
    \item Annoncer que la réponse \NEK{} est admise et documentée, et qu'un
          document non remis sera coté comme tel.
    \item Remettre la liste des documents sollicités (O3) avec une date de
          remise.
    \item Consigner la réunion dans un compte rendu, premier papier de
          travail de la mission.
\end{enumerate}
\end{boxmetier}

% ============================================================================
\section{Le Contradictoire}
\label{sec:c05-contradictoire}
% ============================================================================

\subsection{Pourquoi le contradictoire~?}

\begin{boxdef}
\termecle{Contradictoire}\index{contradictoire}~: principe selon lequel
l'audité prend connaissance des constats qui le concernent et peut y
répondre, avant que le rapport ne soit définitif~; ses réponses sont
consignées et, le cas échéant, reproduites dans le rapport.
\end{boxdef}

Le contradictoire n'est pas une politesse. Il remplit trois fonctions~:
\begin{enumerate}
    \item \textbf{Corriger les erreurs} de l'auditeur, qui a pu mal
          comprendre un processus ou manquer une pièce.
    \item \textbf{Rendre les conclusions opposables}~: un constat discuté et
          maintenu est plus difficile à contester qu'un constat découvert
          dans le rapport (axiome~6).
    \item \textbf{Réduire l'incertitude}~: l'audité apporte souvent, au
          moment du contradictoire, l'information qui manquait.
\end{enumerate}

\subsection{Le contradictoire réduit l'incertitude}

\begin{boxlecture}{Domaine B, avant et après la restitution}
Douze questions du domaine~B (gestion des sauvegardes) sont cotées avant la
réunion de restitution, puis après, lorsque le responsable informatique
remet les journaux de sauvegarde qu'il n'avait pas fournis.

\medskip
\begin{tableaucours}{L{3.8cm}ccccc}{%
    \tableauHeader{Moment} &
    \tableauHeader{\OUI{}} & \tableauHeader{\NON{}} & \tableauHeader{\NEK{}} &
    \tableauHeader{$U$} & \tableauHeader{$H$ (bit)}
}
Avant la restitution & 6 & 2 & 4 & 0,33 & 1,46 \\
Après la restitution & 8 & 3 & 1 & 0,08 & 1,19 \\
\end{tableaucours}

\medskip
\textbf{Constatation.} Le taux d'incertitude passe de 0,33 à 0,08~: trois
\NEK{} sont levés, deux en \OUI{}, un en \NON{}.

\medskip
\textbf{Analyse} (confiance élevée). Le contradictoire a surtout réduit
l'ignorance. L'entropie baisse moins, parce que les réponses restent
partagées entre \OUI{} et \NON{}~: le domaine est désormais mieux connu,
et l'on sait où il est défaillant. Sans restitution, le rapport aurait
conclu sur un tiers de réponses inconnues.

\medskip
\hfill\lecture{0{,}67}{0{,}08}{1{,}19}
\end{boxlecture}

\begin{boxprocedure}
\textbf{La trace du contradictoire.} Pour qu'un constat soit opposable, le
dossier de la mission contient~:
\begin{itemize}
    \item le projet d'observation transmis à l'audité, daté~;
    \item la réponse écrite de l'audité, ou le compte rendu de la réunion de
          restitution signé des participants~;
    \item la position finale de l'auditeur~: constat maintenu, modifié ou
          retiré, et pourquoi~;
    \item les pièces nouvelles remises au moment du contradictoire, avec leur
          date de remise.
\end{itemize}
\end{boxprocedure}

\begin{boxalert}
\textbf{Violation typique.} Accepter, au contradictoire, une affirmation de
l'audité («~les sauvegardes sont testées chaque mois~») et changer un
\NON{} en \OUI{} sans pièce.

\smallskip
\textbf{Règle~:} au contradictoire, une affirmation sans pièce peut lever un
malentendu, jamais remplacer une preuve. Elle transforme au mieux un \NON{}
en \NEK{}, en attendant la pièce.
\end{boxalert}

\begin{boxmetier}{RSSI}
\textbf{Se préparer à être audité}
\begin{enumerate}
    \item Remettre les documents sollicités dans les délais~: un document
          non remis sera coté \NEK{}, et signalé comme tel.
    \item Désigner un interlocuteur unique pour l'équipe d'audit.
    \item Répondre par écrit aux projets d'observation, pièce à l'appui.
    \item Proposer, pour chaque recommandation acceptée, un responsable et
          une échéance.
\end{enumerate}
\end{boxmetier}

\separateur

\begin{boxresume}
\textbf{Ce qu'il faut retenir de ce chapitre~:}
\begin{itemize}
    \item \textbf{Sept principes} guident l'auditeur selon l'ISO~19011:2018,
          dont l'approche fondée sur la preuve et l'approche par les risques.
    \item \textbf{La charte d'audit} fonde les pouvoirs et l'indépendance de
          l'audit~; elle est approuvée par l'organe de gouvernance.
    \item \textbf{Cinq phases}~: préparation, planification, exécution,
          rapport, suivi~; le suivi prépare la mission suivante.
    \item \textbf{Dix outils} (O0 à O9) outillent le cycle~; le classeur O4
          porte la décomposition, la cotation et les mesures.
    \item \textbf{Le contradictoire} corrige, rend opposable et réduit
          l'incertitude~; une affirmation sans pièce ne remplace jamais une
          preuve.
\end{itemize}
\end{boxresume}

\bigskip

\begin{boxexo}[title={\ding{45}\quad Exercice~5.1. Quel principe~? \niveaubase}]
Indiquez le principe de l'ISO~19011 que chacune des situations suivantes
méconnaît.
\begin{enumerate}[(a)]
    \item Un auditeur évoque, lors d'un dîner, les failles du système de
          paiement de la BANQUE SANAGA.
    \item Un rapport passe sous silence un constat qui met en cause le
          directeur général.
    \item Un auditeur consacre autant de temps aux imprimantes qu'aux
          comptes à privilèges.
    \item Un constat repose sur ce qu'un agent a « entendu dire ».
\end{enumerate}
\end{boxexo}

\begin{boxexo}[title={\ding{45}\quad Exercice~5.2. La charte d'audit de LiZbiZ \niveaumoyen}]
LiZbiZ veut créer une fonction d'audit interne.
\begin{enumerate}[(a)]
    \item Rédigez les cinq rubriques principales de sa charte d'audit.
    \item À qui la directrice de l'audit doit-elle être rattachée, et
          pourquoi~?
    \item Quels pouvoirs d'accès faut-il prévoir pour auditer l'application
          intégrée et les quatre centres de données~?
\end{enumerate}
\end{boxexo}

\begin{boxexo}[title={\ding{45}\quad Exercice~5.3. Le contradictoire à la BANQUE SANAGA \niveauexpert}]
Mme~EKAMBI présente à la DSI un projet d'observation~: aucun test de retour
arrière n'a été réalisé avant la bascule de mars~2026. La DSI répond que des
tests ont eu lieu, mais que leurs comptes rendus ont été perdus.
\begin{enumerate}[(a)]
    \item Quelle cotation retenir pour ce contrôle, avant et après la
          réponse de la DSI~? Justifiez.
    \item Que doit contenir le dossier pour que le constat final soit
          opposable~?
    \item Rédigez la position finale de l'auditrice.
    \item En quoi cette situation illustre-t-elle la différence entre réduire
          l'incertitude et lever un constat~?
\end{enumerate}
\end{boxexo}

\bigskip

\begin{boxinfo}
\textbf{Transition.} La mission commence avant le terrain. Le
Chapitre~\ref{chap:planifier} détaille la préparation et la planification~:
termes de référence, documents sollicités, prise de connaissance du système
d'information et plan d'audit.
\end{boxinfo}
